%%%
%%% Radical "⾅"
%%%
\section*{Radical 134: ``⾅''}\addcontentsline{toc}{section}{Radical 134: ⾅}\addcontentsline{loh}{figure}{\#\#\#\# 134: ⾅}

%%%%%%%%%% 舁 %%%%%%%%%%
\subsection*{舁}\addcontentsline{loh}{figure}{舁}

\begin{Entry}{舁}{9}{⾅}
  \begin{Phonetics}{舁}{yu2}
    \definition{v.}{levantar; elevar | aumentar}
  \end{Phonetics}
\end{Entry}

%%%%%%%%%% 舅 %%%%%%%%%%
\subsection*{舅}\addcontentsline{loh}{figure}{舅}

\begin{Entry}{舅}{13}{⾅}
  \begin{Phonetics}{舅}{jiu4}
    \definition*{s.}{Sobrenome: Jiu}
    \definition[个,位,名,些]{s.}{irmão da mãe; tio materno | irmão da esposa; cunhado | Literário: pai do marido; sogro}
  \end{Phonetics}
\end{Entry}

\begin{Entry}{舅舅}{13,13}{⾅,⾅}
  \begin{Phonetics}{舅舅}{jiu4jiu5}[][HSK 7-9]
    \definition[个,位,名]{s.}{tio; irmão da mãe; é assim que você se dirige ao irmão mais velho ou mais novo de sua mãe; você também pode chamá"-lo de 舅}
  \seealsoref{舅}{jiu4}
  \end{Phonetics}
\end{Entry}

%%%%%%%%%% 舆 %%%%%%%%%%
\subsection*{舆}\addcontentsline{loh}{figure}{舆}

\begin{Entry}{舆}{14}{⾅}
  \begin{Phonetics}{舆}{yu2}
    \definition*{s.}{Sobrenome: Yu}
    \definition{adj.}{Obsoleto: público; popular}
    \definition{s.}{Arcaico: carruagem; carroça | liteira; cadeira de arruar | Obsoleto: área; território}
  \end{Phonetics}
\end{Entry}

\begin{Entry}{舆论}{14,6}{⾅,⾔}
  \begin{Phonetics}{舆论}{yu2lun4}[][HSK 7-9]
    \definition{s.}{opinião pública}
  \synonymref{言论}{yan2lun4}
  \synonymref{议论}{yi4lun4}
  \end{Phonetics}
\end{Entry}

%%%%% EOF %%%%%

